\documentclass[11pt,a4paper]{article}

\usepackage[a4paper,margin=2.3cm]{geometry}
\usepackage{fontspec}
\usepackage{titlesec}
\usepackage{longtable}
\usepackage{booktabs}
\usepackage{array}
\usepackage{xcolor}
\usepackage[most]{tcolorbox}
\usepackage{hyperref}
\usepackage{enumitem}
\usepackage{fancyhdr}
\usepackage{amssymb}

\setmainfont{Sylfaen}

\definecolor{brandgreen}{HTML}{1F6F5C}
\definecolor{quoteblue}{HTML}{1D4E89}
\definecolor{termorange}{HTML}{B5651D}
\definecolor{importred}{HTML}{B03A2E}

\hypersetup{
  colorlinks=true,
  linkcolor=brandgreen,
  urlcolor=brandgreen,
  pdftitle={Suttirat et al. 2026 -- ქართული თარგმანი და შენიშვნები},
  pdfauthor={სერგი კონიაშვილი}
}

\titleformat{\section}{\large\bfseries\color{brandgreen}}{\thesection}{0.5em}{}
\titleformat{\subsection}{\normalsize\bfseries}{\thesubsection}{0.5em}{}

\setlength{\headheight}{14pt}
\pagestyle{fancy}
\fancyhf{}
\fancyhead[L]{Suttirat et al. (2026), PLOS NTD}
\fancyhead[R]{ქართული თარგმანი სასწავლად}
\fancyfoot[C]{\thepage}
\setlength{\emergencystretch}{3em}
\setlength{\parskip}{0.4em}

% --- მონიშვნის ყუთები ---
\newtcolorbox{etl}{enhanced,breakable,colback=brandgreen!6,colframe=brandgreen,
  fonttitle=\bfseries,title={ETL: მნიშვნელოვანია ჩემი სისტემისთვის}}
\newtcolorbox{citat}{enhanced,breakable,colback=quoteblue!5,colframe=quoteblue,
  fonttitle=\bfseries,title={ციტატა თეზისისთვის (ორიგინალი ინგლისურად)}}
\newtcolorbox{termin}{enhanced,breakable,colback=termorange!7,colframe=termorange,
  fonttitle=\bfseries,title={ტერმინები}}
\newtcolorbox{skipbox}{enhanced,breakable,colback=gray!7,colframe=gray!55,
  fonttitle=\bfseries,title={ნაკლებად მნიშვნელოვანია შენთვის: საკმარისია ზოგადი გაგება}}
\newcommand{\imp}[1]{\textbf{\color{importred}#1}}
\newcommand{\pg}[1]{{\small\color{gray}[ორიგინალი PDF, გვ.~#1]}}

\begin{document}

\begin{center}
  {\Large\bfseries Machine learning reveals temperature as a key predictor of dengue risk across Thailand's provinces: A 20-year analysis}\\[0.6em]
  {\large მანქანური სწავლება აჩვენებს, რომ ტემპერატურა დენგეს რისკის მთავარი პრედიქტორია ტაილანდის პროვინციებში: 20-წლიანი ანალიზი}\\[0.8em]
  Suttirat P., Chadsuthi S., Aekthong S., Rocklöv J., Bicout D.J., Haddawy P., Su Yin M., Lawpoolsri S., Modchang C.\\
  \textit{PLOS Neglected Tropical Diseases} 20(7): e0014590, 28.07.2026\\
  DOI: \href{https://doi.org/10.1371/journal.pntd.0014590}{10.1371/journal.pntd.0014590}\\
  მონაცემები და კოდი: \href{https://doi.org/10.5281/zenodo.21053822}{10.5281/zenodo.21053822}
\end{center}

\noindent\rule{\linewidth}{0.4pt}

\noindent\textbf{ეს დოკუმენტი სასწავლო თარგმანია და ოფიციალური არ არის.} სტატია ღიაა
(Creative Commons Attribution, CC BY), ამიტომ წყაროს მითითებით თარგმნა დასაშვებია.
თეზისში \imp{ყოველთვის ორიგინალი ინგლისური ტექსტი} ციტირე, ეს თარგმანი კი მხოლოდ გაგებისთვისაა.
Abstract, Introduction, Methods, Limitations და Conclusions სრულადაა თარგმნილი.
Results-სა და Discussion-ში ზოგი აბზაცი შემოკლებულია, ისინი მონიშნულია: \textit{[შემოკლებით]}.

\subsection*{როგორ წაიკითხო ეს დოკუმენტი}
\begin{itemize}[leftmargin=1.5em]
  \item \textcolor{brandgreen}{\textbf{მწვანე ყუთი (ETL)}}: მონაცემთა მომზადების ნაბიჯი, რომელიც ავტორებმა ხელით გააკეთეს. \imp{ეს არის მთავარი}: აქედან გამოიყვან შენი სისტემის მოთხოვნებს.
  \item \textcolor{quoteblue}{\textbf{ლურჯი ყუთი}}: ორიგინალი წინადადება, რომელიც თეზისში შეგიძლია ციტირო.
  \item \textcolor{termorange}{\textbf{ნარინჯისფერი ყუთი}}: ტერმინების ახსნა.
  \item \textcolor{gray}{\textbf{ნაცრისფერი ყუთი}}: ნაწილი, რომელიც შენთვის ნაკლებად მნიშვნელოვანია (მაგ. მოდელის სტატისტიკური დეტალები).
  \item \imp{წითელი ტექსტი}: მნიშვნელოვანი ფაქტი, რომელიც უნდა დაიმახსოვრო.
\end{itemize}

\begin{tcolorbox}[colback=importred!5,colframe=importred,fonttitle=\bfseries,title={ყურადღება: ბრემენის კავშირი}]
ერთ-ერთი თანაავტორი, \textbf{Peter Haddawy}, სტატიაში ასეა მითითებული:
\textit{„Bremen Spatial Cognition Center, University of Bremen, Bremen, Germany“} (გვ.~1).
შენი ხელმძღვანელის წერილი იწყება სიტყვებით \textit{„Hi Peter, hi Lin“}. \imp{შესაძლოა, ეს იგივე პიროვნება იყოს.}
ეს გადაამოწმე. თუ ასეა, ეს სტატია შენი Case Study-სთვის კიდევ უფრო ძლიერი არჩევანია,
რადგან მისი თანაავტორი თავად შეძლებს შეაფასოს, დაეხმარებოდა თუ არა მათ შენი სისტემა.
\end{tcolorbox}

%==========================================================
\section{Abstract (რეზიუმე)} \pg{1--2}

\textbf{ფონი.} დენგეს ცხელება ტაილანდში საზოგადოებრივი ჯანმრთელობის კრიტიკულ პრობლემად რჩება.
მისი გადაცემის დინამიკას გარემოსა და სოციოეკონომიკური ფაქტორების რთული ურთიერთქმედება განსაზღვრავს.
ამ პროგნოზული ფაქტორების გაგება აუცილებელია საიმედო საპროგნოზო სისტემების შესაქმნელად.

\textbf{მეთოდები.} შევქმენით მანქანური სწავლების ჩარჩო (framework), რომელიც დენგეს სივრცით-დროით რისკს
ახარისხებს და მთავარ პროგნოზულ ფაქტორებს ავლენს ტაილანდის მასშტაბით. გავაანალიზეთ
\imp{დენგეს ჰემორაგიული ცხელების (DHF) 20 წლის თვიური ზედამხედველობის მონაცემები (2003--2022) 77 პროვინციიდან}
და მათში \imp{54 გარემოს, კლიმატური და სოციოეკონომიკური ცვლადი} გავაერთიანეთ.
ოთხი კლასიფიკატორი შევადარეთ (ლოგისტიკური რეგრესია, SVM, random forest, XGBoost) და
ექვსი მეტრიკის მიხედვით XGBoost შევარჩიეთ. ცვლადების წვლილის ინტერპრეტაციისთვის SHAP გამოვიყენეთ.
მონაცემები გაიყო სასწავლო (2003--2016) და სატესტო პერიოდებად. ეს უკანასკნელი დაიყო
COVID-19-მდე (2017--2019), COVID-19-ის (2020--2021) და COVID-19-ის შემდგომ (2022) ფაზებად.

\textbf{შედეგები.} XGBoost მოდელმა პანდემიამდე ტესტირებისას AUC = 0.80 მიიღო, პანდემიისა და მის შემდგომ
პერიოდში კი AUC = 0.74. ცვლადების მნიშვნელობის რეიტინგში ტემპერატურა დომინირებდა: ათი უმნიშვნელოვანესი
ცვლადიდან შვიდი ტემპერატურას ეხებოდა. გამოვლინდა არაწრფივი ზღვრები:
\imp{დაახლ. 21°C ერთი თვით დაგვიანებული მინიმალური ტემპერატურისთვის და დაახლ. 32°C სამი თვით დაგვიანებული მაქსიმალური ტემპერატურისთვის}.
ეს მნიშვნელობები \textit{Aedes aegypti}-ს მიერ გადაცემის ცნობილ ბიოლოგიურ შეზღუდვებს ემთხვევა.
ნალექის ცვლადების წვლილი მინიმალური იყო. პროვინციის მაღალი მთლიანი პროდუქტი (GPP)
დენგეს გაზრდილ რისკთან იყო დაკავშირებული, რაც ძირითადად ურბანულ გადაცემას შეესაბამება.
COVID-19-ის დროს მოდელის ხარისხი მნიშვნელოვნად გაუარესდა (2021 წელს AUC = 0.62) და მოდელი
სისტემატურად აჭარბებდა რისკს. ეს მიუთითებს, რომ ამ პერიოდში დენგეს დინამიკას ძირითადად
მოდელის გარეთ არსებული, არაგარემოსეული ფაქტორები განსაზღვრავდა.

\textbf{დასკვნები.} ტემპერატურა დენგეს რისკის დომინანტური პრედიქტორია ტაილანდში, ხოლო აღმოჩენილი
ზღვრები ვექტორის (კოღოს) კომპეტენციის ცნობილ ბიოლოგიურ შეზღუდვებს ემთხვევა. გარემოზე დაფუძნებული
პროგნოზი სტაბილურ პირობებში საიმედოა, მაგრამ დიდი საზოგადოებრივი რყევების დროს არსებითად უარესდება.
ამიტომ, რეალურ დროში გამოყენებისას გარემოს პრედიქტორებს ქცევისა და ზედამხედველობის მოცვის
ინდიკატორებიც უნდა დაემატოს.

\begin{termin}
\textbf{Dengue hemorrhagic fever (DHF)}: დენგეს მძიმე ფორმა. ავტორებმა DHF აიღეს, რადგან მსუბუქი შემთხვევები (DF) ხშირად არ აღირიცხება.\\
\textbf{Classifier / კლასიფიკაცია}: მოდელი, რომელიც კატეგორიას პროგნოზირებს (აქ: „მაღალი“ ან „დაბალი“ რისკი), და არა რიცხვს.\\
\textbf{XGBoost}: გადაწყვეტილების ხეების ანსამბლზე დაფუძნებული ML მეთოდი.\\
\textbf{SHAP}: მეთოდი, რომელიც აჩვენებს, რამდენად შეუწყო ხელი თითოეულმა ცვლადმა მოდელის პროგნოზს.\\
\textbf{AUC}: მოდელის ხარისხის საზომი. 0.5 შემთხვევითი გამოცნობაა, 1.0 იდეალური.\\
\textbf{Lag (დაყოვნება)}: კლიმატის მნიშვნელობა რამდენიმე თვით ადრე. მაგ. \texttt{tmin-1} წინა თვის მინიმალური ტემპერატურაა.
\end{termin}

%==========================================================
\section{Author summary (ავტორების მოკლე შეჯამება)} \pg{2}

ამ კვლევაში მანქანური სწავლება გამოვიყენეთ დენგეს აფეთქებების რისკის საპროგნოზოდ ტაილანდის პროვინციებში
2003--2022 წლებში. \imp{კლიმატის მონაცემების, ეკონომიკური ინდიკატორებისა და სატელიტური გამოსახულებების} ანალიზით
ჩვენი მოდელი მაღალი რისკის რეგიონებს შემთხვევების დაახლოებით 80\%-ში სწორად ავლენდა.
აღმოვაჩინეთ, რომ ტემპერატურა უმნიშვნელოვანესი ფაქტორია, რომელიც განსაზღვრავს, სად და როდის გავრცელდება დენგე.
ასევე გამოვავლინეთ კრიტიკული ზღვარი: როცა მინიმალური ტემპერატურა დაახლოებით 21°C-ზე დაბლა ეცემა,
გადაცემა პრაქტიკულად ჩერდება. საინტერესოა, რომ ნალექს, რომელსაც დენგეს პროგნოზებში ხშირად დიდ
მნიშვნელობას ანიჭებენ, მოსალოდნელზე გაცილებით მცირე როლი ჰქონდა. უფრო მაღალი ეკონომიკური განვითარებაც
დენგეს გაზრდილ რისკთან იყო დაკავშირებული. სავარაუდოდ, ეს ურბანიზაციის გამოა, რომელიც კოღოების
მთელი წლის განმავლობაში გამრავლებისთვის იდეალურ პირობებს ქმნის და ვირუსის გადაცემას მოსახლეობის მაღალი
სიმჭიდროვითაც უწყობს ხელს. თუმცა, COVID-19 პანდემიამ მოდელის სიზუსტე დაარღვია: 2020--2022 წლებში
მოდელი რეალურზე მაღალ რისკს პროგნოზირებდა, სავარაუდოდ COVID-19-ის კონტროლის ზომების გამო,
რომლებსაც მხოლოდ კლიმატის მონაცემები ვერ ასახავს. კვლევა აჩვენებს გარემოსა და სოციოეკონომიკური
მონაცემებით დენგეს რისკის პროგნოზირების როგორც ძალას, ისე შეზღუდვებს.

%==========================================================
\section{Introduction (შესავალი)} \pg{3--4}

დენგეს ცხელება მსოფლიოში ერთ-ერთი ყველაზე სწრაფად გავრცელებადი ვექტორული დაავადებაა: ყოველწლიურად
დაახლოებით 390 მილიონი ინფექცია ხდება 100-ზე მეტ ენდემურ ქვეყანაში, ტროპიკულ და სუბტროპიკულ რეგიონებში [1--3].
ბოლო ხუთი ათწლეულის განმავლობაში დაავადების გლობალური ტვირთი მკვეთრად გაიზარდა, ინციდენტობა 30-ჯერ მოიმატა [4].
ეს ზრდა განსაკუთრებით შესამჩნევია სამხრეთ-აღმოსავლეთ აზიაში, რომელზეც დენგეს გლობალური ტვირთის დაახლოებით
70\% მოდის [1]. სამხრეთ-აღმოსავლეთ აზიის ქვეყნებს შორის ტაილანდი მუდმივად რეგიონის ერთ-ერთ ყველაზე მაღალ
ინციდენტობას აფიქსირებს, რაც მას დენგეს გადაცემის დინამიკის შესწავლის კრიტიკულ ცენტრად აქცევს [5].

დენგეს ვირუსს ოთხი განსხვავებული სეროტიპი აქვს და ადამიანს ძირითადად ინფიცირებული \textit{Aedes} კოღოს
კბენით გადაეცემა [2,3,6]. კლინიკური სურათი ძალიან განსხვავდება: უსიმპტომო ინფექციიდან ან მსუბუქი
ცხელებიდან კლასიკურ დენგეს ცხელებამდე (DF), მძიმე შემთხვევებში კი სიცოცხლისთვის საშიშ მდგომარეობამდე,
როგორიცაა დენგეს ჰემორაგიული ცხელება (DHF) ან დენგეს შოკის სინდრომი (DSS) [2,3]. ყველაზე მძიმე
ფორმებმა შეიძლება პლაზმის გაჟონვა, სისხლდენა და ორგანოების დაზიანება გამოიწვიოს [7]. ტაილანდმა
DHF-ის პირველი ეპიდემია 1958 წელს დააფიქსირა [8,9]. მას შემდეგ დაავადება საზოგადოებრივი ჯანმრთელობის
მუდმივ გამოწვევად იქცა, რომელსაც კლიმატური პირობების, სოციოეკონომიკური ფაქტორებისა და ვექტორის ეკოლოგიის
ურთიერთქმედებით განპირობებული რთული სივრცით-დროითი სურათი აქვს [10,11].

დენგეს ეპიდემიოლოგიას სოციოეკონომიკური, დემოგრაფიული და ინფრასტრუქტურული ფაქტორების რთული
ურთიერთქმედება განსაზღვრავს [12--14]. რამდენიმე კვლევამ დაბალი სოციოეკონომიკური სტატუსი გაზრდილ
ადგილობრივ რისკთან დააკავშირა [13,15], თუმცა კავშირი უფრო რთულია. სწრაფი ურბანიზაცია, მოსახლეობის
მაღალი სიმჭიდროვე და ნარჩენების არასათანადო მართვა დაავადების გავრცელებისთვის ხელსაყრელ პირობებს ქმნის [16,17].
თუმცა უახლესი მონაცემები ეჭვქვეშ აყენებს ტრადიციულ შეხედულებას, რომ დენგე სიღარიბის ურბანული დაავადებაა.
ვირუსი ახლა სოფლისა და ქალაქის მოსახლეობას სხვადასხვა სოციოეკონომიკურ ფენაში აზიანებს [18,19].

გარემო ფაქტორებს დენგეს გადაცემაში გადამწყვეტი როლი აქვს. \imp{კლიმატური ცვლადები, განსაკუთრებით
ტემპერატურა და ნალექი}, ურბანული განვითარების ხარისხთან ერთად, ადგილობრივი რისკის პროფილს
ფუნდამენტურად განსაზღვრავს [1,20,21]. ტემპერატურა გადაცემაზე განსაკუთრებით ძლიერად მოქმედებს,
რადგან კოღოს ბიოლოგიურ პარამეტრებს ცვლის: კბენის სიხშირეს, ნაყოფიერებას, განვითარების სიჩქარეს,
ზრდასრული კოღოს გადარჩენას [22,23].

\begin{citat}
\textit{„Traditional approaches to dengue forecasting have relied predominantly on epidemiological surveillance data combined with basic statistical models that incorporate rainfall and temperature variables [24--26].“} (Introduction, p.~3)
\end{citat}

დენგეს პროგნოზირების ტრადიციული მიდგომები ძირითადად ეპიდემიოლოგიური ზედამხედველობის მონაცემებს
ეყრდნობოდა, რომლებიც ნალექისა და ტემპერატურის ცვლადებთან ერთად მარტივ სტატისტიკურ მოდელებში
გამოიყენებოდა [24--26]. თუმცა ეს მეთოდები ხშირად ვერ ასახავს რთულ, არაწრფივ კავშირებს მრავალ გარემო
და სოციოეკონომიკურ ფაქტორს შორის [27]. ამან მკვლევრები უფრო დახვეწილი მეთოდების ძიებისკენ უბიძგა.

მანქანური სწავლების ბოლოდროინდელი მიღწევები დაავადებების პროგნოზირების პერსპექტიულ ალტერნატივებს
გვთავაზობს. ანსამბლური მეთოდები, მაგალითად gradient boosting, ცვლადებს შორის რთული ურთიერთქმედებების
აღმოჩენაში შესანიშნავ შედეგებს აჩვენებს [27--29]. ამ მეთოდებს \imp{მაღალგანზომილებიანი მონაცემების}
დამუშავება შეუძლია.

\begin{citat}
\textit{„Despite growing interest in applying machine learning to dengue prediction, significant knowledge gaps still persist [28,30]. Most existing models suffer from a limited geographic scope or analyze only short time periods [28,30,31].“} (Introduction, p.~3)
\end{citat}

მიუხედავად დენგეს პროგნოზირებაში მანქანური სწავლების მიმართ მზარდი ინტერესისა, ცოდნაში მნიშვნელოვანი
ხარვეზები რჩება [28,30]. \imp{არსებული მოდელების უმეტესობას შეზღუდული გეოგრაფიული მასშტაბი აქვს ან
მხოლოდ მოკლე დროის პერიოდს აანალიზებს} [28,30,31]. ბევრი მათგანი ნალექს მთავარ პრედიქტორად
არაპროპორციულად ეყრდნობა, მიუხედავად იმისა, რომ ნალექსა და დენგეს შორის კავშირი არათანმიმდევრულია [32].
ამასთან, მხოლოდ ცოტა მოდელი ხსნის, რომელი ცვლადები განაპირობებს მის პროგნოზს, რაც საზოგადოებრივი
ჯანდაცვის გადაწყვეტილებებისთვის მათ სარგებელს ზღუდავს [27].

\begin{etl}
ეს შენი მოტივაციისთვის ძალიან მნიშვნელოვანია: ავტორები ამბობენ, რომ \imp{მოდელების უმეტესობა მცირე
გეოგრაფიულ არეალს ან მოკლე პერიოდს მოიცავს}. ერთ-ერთი შესაძლო მიზეზი ისაა, რომ დიდი, მრავალწლიანი
და მრავალწყაროიანი მონაცემთა ნაკრების მომზადება შრომატევადია. \textbf{მაგრამ ავტორები ამ მიზეზს პირდაპირ არ ასახელებენ!}
თეზისში ასე არ დაწერო: „Suttirat-ი ამბობს, რომ მონაცემების მომზადება რთულია“. დაწერე ის, რაც ნამდვილად წერია
(მოდელები შეზღუდულია მასშტაბით), ხოლო საკუთარი ინტერპრეტაცია ცალკე, როგორც შენი არგუმენტი.
\end{etl}

ამასთან, ცუდად არის შესწავლილი, როგორ მოქმედებს საზოგადოებრივი რყევები გარემოზე დაფუძნებული
მოდელების ხარისხზე. ძალიან ცოტა კვლევა აფასებს სისტემატურად, როგორ ცვლის ამას ისეთი მოვლენა, როგორიც
COVID-19 პანდემიაა [33]. პანდემიის ზომები, მათ შორის მობილობის შეზღუდვები და ლოკდაუნი, სავარაუდოდ
დენგეს გადაცემას გარემო პირობებისგან დამოუკიდებელი მექანიზმებით ცვლიდა [33--36].

ამ ხარვეზების აღმოსაფხვრელად ეს კვლევა ქმნის და აფასებს მანქანური სწავლების ჩარჩოს დენგეს სივრცით-დროითი
რისკის კლასიფიკაციისთვის ტაილანდში. ვიყენებთ XGBoost-ს SHAP-თან ერთად, რათა გავაანალიზოთ \imp{ორი ათწლეულის
ზედამხედველობის მონაცემები (2003--2022), გაერთიანებული გარემოსა და სოციოეკონომიკურ ცვლადებთან}.
DHF შემთხვევებზე გავამახვილეთ ყურადღება, რადგან ისინი კლინიკურ სიმძიმესა და საიმედო აღრიცხვას შორის
ბალანსს იძლევა. მსუბუქი DF შემთხვევები სისტემატურად არ აღირიცხება [37--39], ხოლო DSS ძალიან იშვიათია
პროვინციის დონის თვიური კლასიფიკაციისთვის. 18\,240 პროვინცია-თვიდან 59\%-ში DSS საერთოდ არ
დაფიქსირებულა, ხოლო DHF-ის მედიანური ინციდენტობა (2.78 / 100\,000) მაღალი და დაბალი რისკის
დაბალანსებულ გაყოფას იძლევა. კონკრეტული მიზნები სამია:
(1) მოდელის შექმნა, რომელიც სხვადასხვა გარემო და სოციოეკონომიკური ფაქტორის მნიშვნელობასა და
ურთიერთქმედებას შეაფასებს; (2) კრიტიკული ზღვრების აღმოჩენა მთავარ პრედიქტორებში; (3) მოდელის ხარისხის
შეფასება სხვადასხვა ეპიდემიოლოგიურ პერიოდში, განსაკუთრებით COVID-19-ის გავლენის თვალსაზრისით.

\begin{termin}
\textbf{პროვინცია-თვე (province-month)}: მონაცემთა ერთი სტრიქონი: ერთი პროვინცია ერთ კონკრეტულ თვეში. 77 პროვინცია $\times$ 240 თვე = 18\,480. სტატიაში წერია 18\,240, რადგან ერთი პროვინცია მეორესთან გაერთიანდა (იხ. ქვემოთ).\\
\textbf{ინციდენტობა 100\,000-ზე}: შემთხვევების რაოდენობა / მოსახლეობა $\times$ 100\,000. ეს შენს FR-17-ს უკავშირდება.
\end{termin}

%==========================================================
\section{Materials and methods (მასალა და მეთოდები)} \pg{4--8}

\imp{ეს თავი ყველაზე ყურადღებით წაიკითხე, შეიძლება ორჯერაც.} სწორედ აქ ჩანს, რას აკეთებს მეცნიერი
მოდელირებამდე და, შესაბამისად, რისი ავტომატიზება შეუძლია შენს სისტემას.

\subsection*{ეთიკური განცხადება}
კვლევა ეთიკური განხილვისგან გათავისუფლდა (მაჰიდოლის უნივერსიტეტი), რადგან ანონიმიზებული,
\imp{აგრეგირებული} ზედამხედველობის მონაცემების რეტროსპექტულ ანალიზს ეყრდნობოდა.

\subsection*{მეთოდოლოგიის სამუშაო ნაკადი (Methodology workflow)} \pg{4}

ამ კვლევამ შექმნა ჩარჩო დენგეს რისკის პრედიქტორების გამოსავლენად (S1 Fig). ჩარჩო მანქანური სწავლებით
ყველაზე ინფორმაციულ პრედიქტორებს ახარისხებს \imp{54 შემავალ ცვლადს შორის: კლიმატური, ბიოკლიმატური,
სოციოეკონომიკური და დისტანციური ზონდირების (სატელიტური) მონაცემები}. მონაცემები \imp{სამი ძირითადი წყაროდან}
იქნა მიღებული: TerraClimate, MODIS და ტაილანდის ეროვნული უწყებები. \imp{გრიდული კლიმატური და სატელიტური
მონაცემები სივრცით აგრეგირდა პროვინციის დონეზე; შემდეგ პროვინციების თვიური ტემპერატურისა და ნალექის
სერიებიდან გამოითვალა დაყოვნებული (lagged) ცვლადები და ბიოკლიმატური ინდექსები.} სოციოეკონომიკური
მონაცემები პირდაპირ პროვინციის დონეზე იქნა მიღებული: \imp{ორ წელიწადში ერთხელ გამოქვეყნებული შინამეურნეობის
შემოსავლის სერია წრფივად ინტერპოლირდა} გამოტოვებული წლების შესავსებად, შემდეგ კი \imp{ყველა სოციოეკონომიკური
სერია თვიურ რეზოლუციაზე გადავიდა: ყოველი წლის მნიშვნელობა ამ წლის ყველა თვეს მიენიჭა}. თვიური ცვლადების
ნაკრების აგების შემდეგ ოთხი ML მოდელი შედარდა წინასწარ (S1 Table). შეირჩა XGBoost; ჰიპერპარამეტრები
ბაიესური ძიებით და 5-ჯერადი ჯვარედინი ვალიდაციით (cross-validation) იქნა ოპტიმიზებული.

\begin{etl}
ეს აბზაცი \textbf{სტატიის მთელ მონაცემთა მილსადენს (pipeline)} აღწერს. ნაბიჯ-ნაბიჯ:
\begin{enumerate}[leftmargin=1.5em]
  \item მონაცემთა მოპოვება \textbf{სამი} განსხვავებული წყაროდან (კლიმატი, სატელიტი, სახელმწიფო სტატისტიკა);
  \item გრიდის \textbf{სივრცითი აგრეგაცია} პროვინციებზე;
  \item \textbf{lag ცვლადების} გამოთვლა;
  \item \textbf{ბიოკლიმატური ინდექსების} გამოთვლა;
  \item გამოტოვებული წლების \textbf{წრფივი ინტერპოლაცია};
  \item წლიური მონაცემების \textbf{თვიურზე გადაყვანა (upsampling)};
  \item ყველაფრის \textbf{ერთ თვიურ ცხრილად} გაერთიანება.
\end{enumerate}
ეს ზუსტად \textbf{ETL}-ია (Extract, Transform, Load). \imp{ეს სია შენი თეზისის მოტივაციის ბირთვია:}
ყოველ ნაბიჯთან ჩაწერე, შეუძლია თუ არა ეს შენს სისტემას.
\end{etl}

\begin{citat}
\textit{„Gridded climatic and remote-sensing data were spatially aggregated to the provincial level; lagged variables and bioclimatic indices were then derived from monthly provincial temperature and precipitation series. Socioeconomic data were obtained directly at the provincial level: the biennially reported household income series was linearly interpolated to fill missing years, and all socioeconomic series were then upsampled to monthly resolution by assigning each year's reported value (typically released at year-end) to every month of that year.“} (Methods > Methodology workflow, p.~4)
\end{citat}

\subsection*{კვლევის დიზაინი და მონაცემთა წყაროები} \pg{5}

\textbf{დენგეს შემთხვევების მონაცემები.} კვლევაში გამოყენებულია \imp{DHF შემთხვევების თვიური ანგარიშები,
მიღებული ტაილანდის ჯანდაცვის სამინისტროს დაავადებათა კონტროლის დეპარტამენტის ეპიდემიოლოგიის ბიუროდან},
2003 წლის იანვრიდან 2022 წლის დეკემბრამდე. შემთხვევის განსაზღვრება ეროვნულ სტანდარტებს მიჰყვებოდა:
დადებითი შემთხვევა კლინიკური შეფასებით ან ლაბორატორიული დადასტურებით დგინდებოდა. პროვინციების მონაცემები
თვეების მიხედვით აგრეგირდა. \imp{ბუენგკანის (Bueng Kan) პროვინციის შემთხვევები ნონგ ხაის (Nong Khai)
პროვინციას მიემატა, რადგან ბუენგკანი 2011 წელს ადმინისტრაციულად გამოეყო ნონგ ხაის. ეს თანმიმდევრულობისთვის
გაკეთდა.} ეპიდემიოლოგიური ერთეულია პროვინცია, დროის ბიჯი კი თვე.

\begin{etl}
\textbf{ადმინისტრაციული საზღვრების ცვლილება დროში} ტიპური და რთული პრობლემაა: 2011 წლამდე ბუენგკანი
არ არსებობდა, ამიტომ 20-წლიან სერიაში ის ხელით უნდა გაერთიანდეს. \imp{შენს სისტემას ეს ჯერ არ შეუძლია}
(ის ქვეყნის დონეზე მუშაობს). ეს შესაძლო მოთხოვნაა („admin-unit harmonisation“) ან შეზღუდვა.
\end{etl}

\textbf{კლიმატური ცვლადები.} კლიმატის მონაცემები ამოღებულია \imp{TerraClimate} ნაკრებიდან [40], რომელიც
მაღალი რეზოლუციის (4 კმ\textsuperscript{2}) გრიდულ კლიმატურ ინფორმაციას იძლევა. მივიღეთ
\imp{ექვსი ძირითადი ცვლადი: ნალექი, მინიმალური ტემპერატურა, მაქსიმალური ტემპერატურა, ნიადაგის ტენიანობა,
ორთქლის წნევა და ქარის სიჩქარე}. \imp{ეს გრიდული მონაცემები პროვინციის დონეზე სივრცით გასაშუალოვდა},
რათა ეპიდემიოლოგიურ ერთეულებს დამთხვეოდა. ამას გარდა, R-ის პაკეტით \texttt{dismo} თვიური ტემპერატურისა და
ნალექის მონაცემებიდან \imp{19 ბიოკლიმატური ინდიკატორი} გამოვთვალეთ [41,42]. ისინი წლიურ ტენდენციებს,
სეზონურობასა და ცვალებადობას ასახავს.

\begin{etl}
\begin{itemize}[leftmargin=1.5em]
  \item \textbf{TerraClimate, არა ERA5!} (Gemini-მ ეს შეცდომით დაწერა.) შენი სისტემა Open-Meteo ERA5-სა და NASA POWER-ს იყენებს. ეს განსხვავება თეზისში გულწრფელად უნდა ახსენო.
  \item \textbf{6 ცვლადი}: შენს სისტემას ახლა 2 აქვს (საშუალო ტემპერატურა და ნალექი). \imp{Tmin და Tmax აკლია}, სტატიაში კი ზუსტად ესენი აღმოჩნდა ყველაზე მნიშვნელოვანი.
  \item \textbf{„spatially averaged to the provincial level“}: შენი სისტემა კლიმატს \imp{ერთი წერტილიდან} იღებს (ერთი lat/lon ქვეყანაზე). სტატია \textbf{ფართობზე გასაშუალოებას} იყენებს. აღსანიშნავია, რომ ავტორები არ წერენ, \textit{როგორ} გაასაშუალოვეს (შეწონილად თუ არა). ესეც კარგი არგუმენტია: დოკუმენტაცია (provenance) არასრულია, შენი სისტემა კი ყოველ ნაბიჯს ავტომატურად აღწერს.
\end{itemize}
\end{etl}

\textbf{სოციოეკონომიკური ინდიკატორები.} ურბანიზაციისა და ეკონომიკური განვითარების ეფექტის გასათვალისწინებლად
სამი სოციოეკონომიკური ფაქტორი ჩავრთეთ. \imp{პროვინციის მთლიანი პროდუქტი (GPP)}, ინფლაციით
კორექტირებული (Laspeyres-ის ინდექსით), მიღებულია ეროვნული ეკონომიკური და სოციალური განვითარების
საბჭოდან (NESDC). \imp{დემოგრაფიული მონაცემები}, მათ შორის პროვინციის მოსახლეობა და შინამეურნეობის საშუალო
ზომა, მიღებულია რეგისტრაციის ადმინისტრაციის ბიუროდან (BORA). პროვინციის შინამეურნეობის საშუალო შემოსავალი
მიღებულია ტაილანდის სტატისტიკის ეროვნული ოფისიდან (NSO), სადაც მნიშვნელობები ორ წელიწადში ერთხელ
ქვეყნდება. \imp{გამოტოვებული წლები წრფივად ინტერპოლირდა pandas-ის \texttt{interpolate} მეთოდით.}
\imp{ყველა სოციოეკონომიკური სერია წლიურიდან თვიურ რეზოლუციაზე გადავიდა: ყოველი წლის მნიშვნელობა ამ წლის
ყველა თვეს მიენიჭა}, იმ დაშვებით, რომ წლიური მნიშვნელობა მთელი წლის პირობებს ასახავს.

\begin{etl}
\begin{itemize}[leftmargin=1.5em]
  \item \textbf{სამი სხვადასხვა სახელმწიფო უწყება} სამი ცალკე ვებგვერდით: მკვლევარმა თითოეული ხელით უნდა მოძებნოს, ჩამოტვირთოს და გაასუფთაოს.
  \item \textbf{მოსახლეობა} ინციდენტობის გამოსათვლელად სჭირდებათ. შენთან FR-17 (World Bank / UN WPP) ამას ქვეყნის დონეზე აკეთებს.
  \item \textbf{წლიური $
\rightarrow$ თვიური}: ზუსტად ის, რასაც შენი სისტემა რეზოლუციების გაერთიანებისას აკეთებს. შეადარე შენს მეთოდს: შენ \textit{mass-preserving} downscaling გაქვს (FR-20), ისინი კი წლის მნიშვნელობას უბრალოდ იმეორებენ ყოველ თვეში. ორივე მიდგომა ლეგიტიმურია, მაგრამ \textbf{სხვადასხვა ტიპის ცვლადისთვის}: შემთხვევების რაოდენობა „მასაა“ (ჯამი უნდა შენარჩუნდეს), შემოსავალი კი „დონე“ (ის უნდა განმეორდეს). კარგი თემაა თეზისისთვის.
  \item \textbf{pandas interpolate}: ავტორები Python-სა და pandas-ს იყენებენ, ანუ იგივე ხელსაწყოებს, რაც შენ.
\end{itemize}
\end{etl}

\textbf{დისტანციური ზონდირების მონაცემები.} გარემო პირობები შეფასდა Terra თანამგზავრის MODIS მონაცემებით
(MOD09A1: ზედაპირის არეკვლა, 8-დღიანი, 500 მ გრიდი). გამოვთვალეთ ორი ეკოლოგიური ინდექსი, MNDWI და NDVI,
ჩვენი მანამდე შემუშავებული მეთოდოლოგიით [43]. წყლის ინდექსი (MNDWI) გამოხატულია, როგორც პროვინციის
პიქსელების პროცენტი, სადაც MNDWI $>$ 0. მცენარეულობის ინდექსი (NDVI) გამოხატულია, როგორც პიქსელების
პროცენტი, სადაც NDVI 0.5--1.0 შუალედშია. ცვლადების სრული აღწერა და წყაროები Table 1-შია (გვ.~6--7).

\begin{termin}
\textbf{NDVI}: მცენარეულობის ინდექსი სატელიტური გამოსახულებიდან. რაც უფრო მწვანეა ტერიტორია, მით მაღალია.\\
\textbf{MNDWI}: წყლის ინდექსი. აჩვენებს, ტერიტორიის რა ნაწილი არის წყლით დაფარული.\\
\textbf{ბიოკლიმატური ინდიკატორები (bio1--bio19)}: სტანდარტული წლიური მაჩვენებლები, მაგ. bio1 = წლის საშუალო ტემპერატურა, bio12 = წლის ჯამური ნალექი.
\end{termin}

\begin{skipbox}
სატელიტური ინდექსები (NDVI/MNDWI) შენი სისტემის ფარგლებს სცილდება. თეზისში შეგიძლია ჩაწერო, რომ
ეს არის „Future Work“ ან შეზღუდვა. ავტორების შედეგებით მათი წვლილი მინიმალური იყო.
\end{skipbox}

\subsection*{Table 1 (გვ.~6--7): 54 ცვლადის ცხრილი}

\imp{ეს ცხრილი ორიგინალში აუცილებლად ნახე.} მოკლე შინაარსი:

\begin{longtable}{@{}p{0.08\linewidth}p{0.30\linewidth}p{0.17\linewidth}p{0.38\linewidth}@{}}
\toprule
\textbf{\#} & \textbf{ცვლადი} & \textbf{რეზოლუცია} & \textbf{წყარო / შენიშვნა}\\
\midrule
1--6 & ppt, tmin, tmax, vap, soil, ws & თვიური & TerraClimate\\
7--24 & ზემოთ ჩამოთვლილი ექვსივე, lag 1, 2, 3 & თვიური & გამოთვლილი (6 $\times$ 3 = 18 სვეტი)\\
25--43 & bio1--bio19 & წლიური & გამოთვლილი \texttt{dismo}-თი\\
44--46 & GPP, შინამეურნეობის ზომა, შემოსავალი & წლიური & NESDC, BORA, NSO\\
47--54 & per.mndwi, per.ndvi + lag 1, 2, 3 & თვიური & MODIS\\
სამიზნე & დენგეს რისკი (0/1) & თვიური & ეპიდ. ბიურო: ინციდენტობა 100k-ზე $\geq$ მედიანა $
\rightarrow$ 1\\
\bottomrule
\end{longtable}

\begin{etl}
შენიშნე, რომ \imp{54 ცვლადიდან 24 lag-ცვლადია} (18 კლიმატური + 6 სატელიტური). ანუ lag-ები მოდელისთვის
\textbf{ცხრილის ცალკე სვეტებად} გამოითვლება. შენი სისტემა lag-ებს ახლა მხოლოდ კორელაციის ანგარიშში
იყენებს (FR-19). ამიტომ Gemini-ს დავალება B (lag სვეტები) ზუსტად ამ ცხრილიდან გამომდინარეობს.
\end{etl}

\subsection*{მონაცემთა დამუშავება და ცვლადების კონსტრუირება (feature engineering)} \pg{5--7}

\textbf{დროითი დაყოფა.} მონაცემები სასწავლო (2003--2016, 14 წელი) და სატესტო (2017--2022, 6 წელი) პერიოდებად
გაიყო. სასწავლო პერიოდი საკმარისია წლიდან წლამდე ცვალებადობისა და რამდენიმე ეპიდემიური ციკლის
დასაჭერად, ბოლო წლები კი დამოუკიდებელი ვალიდაციისთვის დარჩა. სასწავლო ნაწილის ბოლო ორი წელი
(2015--2016) ადრეული შეჩერების (early stopping) ვალიდაციისთვის გამოიყენეს. სატესტო პერიოდი სამ ფაზად დაიყო:
COVID-19-მდე (2017--2019), COVID-19 (2020--2021) და COVID-19-ის შემდეგ (2022).

\textbf{რისკის კლასიფიკაცია.} მიზანი იყო ოპერატიულად გამოსადეგი ადრეული გაფრთხილების სისტემის შექმნა,
რომელიც პროვინცია-თვეებს მაღალი ან დაბალი რისკის კატეგორიად მონიშნავს და შემთხვევების ზუსტ რაოდენობას
არ პროგნოზირებს. ბინარული კლასიფიკაცია საზოგადოებრივი ჯანდაცვისთვის მკაფიო გადაწყვეტილების ზღვარს იძლევა
და სისტემატური აღრიცხვის მიკერძოებისადმი უფრო მდგრადია.

\imp{ყოველი პროვინციისთვის თვიური DHF ინციდენტობა 100\,000 მოსახლეზე გამოვთვალეთ.} სასწავლო პერიოდის
ეროვნული მედიანა (2.78 / 100\,000) კლასიფიკაციის ზღვრად ავიღეთ: ყოველი პროვინცია-თვე მაღალი რისკის
($\geq$ მედიანა) ან დაბალი რისკის ($<$ მედიანა) კატეგორიად მოინიშნა. ეს დაახლოებით 50/50 დაბალანსებულ
კლასებს იძლევა.

\textbf{დაყოვნებული ცვლადები (Lagged variables).} \pg{7}
\imp{გარემოს დაგვიანებული ეფექტის დასაჭერად ყველა კლიმატური და სატელიტური ცვლადისთვის 1, 2 და 3 თვით
დაყოვნებული მნიშვნელობები შევქმენით}, რითაც ცვლადების ნაკრები გაფართოვდა.

\begin{citat}
\textit{„To capture delayed environmental effects on dengue transmission, we generated 1-, 2-, and 3-month lagged values for all climate and remote sensing variables, expanding our feature set to account for temporal dependencies in disease dynamics.“} (Methods > Data processing and feature engineering, p.~7)
\end{citat}

\begin{etl}
\textbf{ინციდენტობა 100k-ზე} = შენი FR-17. \textbf{Lag 1--3} = შენი FR-19 (კორელაცია) და მომავალი lag სვეტები.
ეს ორი ციტატა თეზისში შეგიძლია გამოიყენო, როგორც \imp{პირდაპირი დასაბუთება}: „FR-17 და FR-xx
გამომდინარეობს Suttirat et al. (2026)-დან, სადაც ...“.
\end{etl}

\subsection*{მანქანური სწავლების მიდგომა} \pg{7--8}

\begin{skipbox}
ეს ნაწილი მოდელს ეხება და მონაცემებს არა. შენი სისტემა მოდელს არ აგებს, მხოლოდ მის შემავალ მონაცემებს ამზადებს.
საკმარისია ზოგადად იცოდე:
\begin{itemize}[leftmargin=1.5em]
  \item \textbf{XGBoost}: გადაწყვეტილების ხეები, რომლებიც ერთმანეთის შეცდომებს თანმიმდევრულად ასწორებს.
        ცვლადების სტანდარტიზაცია (მასშტაბირება) არ სჭირდება, ამიტომ ავტორებს ეს არ გაუკეთებიათ.
        ტრენდი და სეზონურობაც არ მოუშორებიათ.
  \item \textbf{SHAP}: ყოველი ცვლადის წვლილი ყოველ პროგნოზში, თამაშთა თეორიაზე დაფუძნებით.
  \item \textbf{Optuna / TPE}: ჰიპერპარამეტრების ავტომატური ოპტიმიზაცია; 5-ჯერადი stratified cross-validation; early stopping 30 რაუნდზე.
  \item \textbf{მეტრიკები}: AUC, Accuracy, Precision, Sensitivity, F1-score (ფორმულები გვ.~8-ზეა).
  \item \textbf{დანერგვა}: Python 3.10, \imp{pandas} (მონაცემთა მანიპულაცია), numpy, scikit-learn, XGBoost, SHAP.
\end{itemize}
\end{skipbox}

\begin{etl}
ერთი წინადადება აქედან მაინც მნიშვნელოვანია: \textit{„pandas for data manipulation“}. ანუ მონაცემთა
მომზადება ავტორებმა საკუთარი Python სკრიპტებით გააკეთეს. შენი სისტემა მომხმარებელს
\imp{პროგრამირების გარეშე} იგივეს აძლევს (NFR: usability).
\end{etl}

%==========================================================
\section{Results (შედეგები)} \pg{9--14}

\begin{skipbox}
შედეგები ეპიდემიოლოგიურ აღმოჩენებს ეხება. შენთვის მნიშვნელოვანია მხოლოდ ის, რომ \textbf{lag-ებიანი
ტემპერატურის ცვლადები} აღმოჩნდა მთავარი. ეს ასაბუთებს, რომ Tmin, Tmax და lag-ები მონაცემთა ნაკრებში აუცილებელია.
\end{skipbox}

\textbf{დენგეს სივრცით-დროითი სურათი} \textit{[შემოკლებით]}. 2003--2022 წლებში ტაილანდში 681\,430 DHF შემთხვევა
დაფიქსირდა. თვიური მაქსიმუმი 13\,601 იყო (2013 წლის ივლისი), მინიმუმი 105 (2022 წლის თებერვალი).
პიკი ადრეულ წვიმიან სეზონზე მოდის: ივნისი (13.16\%), ივლისი (14.96\%) და აგვისტო (13.98\%), სულ 42.12\%.
ყველაზე დაბალი თებერვალშია (3.71\%). ყველაზე მაღალი რეგიონული ინციდენტობა სამხრეთ პროვინციებშია
(13.93 / 100\,000, 2010 წ.). მაღალი რისკის ზონები ცენტრალურ და სამხრეთ ტაილანდშია.

\begin{etl}
Fig.~1-ის რუკები ავტორებმა შექმნეს \imp{HDX}-იდან (Humanitarian Data Exchange) მიღებული ტაილანდის
ადმინისტრაციული საზღვრების shapefile-ით (\texttt{cod-ab-tha}). შენს სისტემას უკვე აქვს \texttt{hdx.py} მოდული!
ეს კარგი საწყისია სუბნაციონალური (პროვინციის დონის) ფუნქციონალისთვის.
\end{etl}

\textbf{მოდელის ხარისხი} \textit{[შემოკლებით]}. XGBoost-ის AUC სასწავლო მონაცემებზე 0.94 იყო, COVID-19-მდე
ტესტზე 0.80, პანდემიის დროს და მის შემდეგ კი 0.74. ნარჩენების (residuals) ავტოკორელაცია მოკლე lag-ებზე
მნიშვნელოვანი აღმოჩნდა (lag-1 = 0.54, lag-2 = 0.37, lag-3 = 0.24), ანუ მოდელი თვიდან თვემდე დამოკიდებულებას
ბოლომდე ვერ ხსნის. წინა თვის რისკის სტატუსის ცვლადად დამატებით ეს პრობლემა მოგვარდა, ტემპერატურა კი
მთავარ ფაქტორად დარჩა. მთლიანი სიზუსტე 20 წელზე 79.27\% იყო.

\textbf{ცვლადების მნიშვნელობა} \textit{[შემოკლებით]}. 54 ცვლადიდან ტოპ-10-ში შვიდი ტემპერატურისაა.
\imp{ორი ყველაზე ძლიერი: ერთი თვით დაყოვნებული მინიმალური ტემპერატურა (tmin-1) და მიმდინარე თვის
მინიმალური ტემპერატურა (tmin). მესამეა სამი თვით დაყოვნებული მაქსიმალური ტემპერატურა (tmax-3).}
მეოთხეა GPP. ტოპ-10 ერთად მოდელის პროგნოზული უნარის 58.42\%-ს შეადგენს. ნალექის, ნალექის
ბიოკლიმატური ინდექსებისა და სატელიტური ცვლადების წვლილი მინიმალური იყო.

SHAP-ის დამოკიდებულების გრაფიკებმა (Fig.~5) ზღვრული ეფექტები აჩვენა: tmin-1 დაახლ. 21°C-ზე დაბლა
რისკს ამცირებს, დადებითი ეფექტი კი დაახლ. 24°C-ზე პლატოს აღწევს. tmax-3-ს კრიტიკული ზღვარი დაახლ.
32°C-ზე აქვს, რის ზემოთაც რისკი მკვეთრად იზრდება.

\textbf{COVID-19-ის გავლენა} \textit{[შემოკლებით]}. 2021 წელს AUC 0.62-მდე დაეცა. 0.5 ზღვრით მოდელმა
2020, 2021 და 2022 წლებისთვის 545, 464 და 581 მაღალი რისკის პროვინცია-თვე იწინასწარმეტყველა,
რეალურად კი მხოლოდ 260, 7 და 121 იყო. ეს მიუთითებს, რომ პანდემიის ფაქტორებმა (მობილობის შეზღუდვა,
ქცევის ცვლილება, ექიმთან მიმართვის შემცირება) დენგეს გადაცემა კლიმატისგან დამოუკიდებლად შეამცირა.

%==========================================================
\section{Discussion (განხილვა)} \pg{14--18}

\textit{[შემოკლებით]}. კვლევამ აჩვენა, რომ XGBoost-ით შესაძლებელია DHF რისკის კლასიფიკაცია
\imp{კლიმატური, გარემოსა და სოციოეკონომიკური ცვლადების ინტეგრაციით}. მოდელის ჩავარდნა პანდემიის დროს
მიუთითებს ფაქტორებზე, რომლებიც მოდელში არ შედიოდა: ადამიანების მობილობა, კონტროლის ზომები და
ჯანდაცვაზე მიმართვის ქცევა.

\textbf{ტემპერატურა, როგორც მთავარი პრედიქტორი} \textit{[შემოკლებით]}. ~21°C ზღვარი
\textit{Aedes aegypti}-ს ბიოლოგიურ შეზღუდვებს ემთხვევა: ამ ტემპერატურის ქვემოთ კბენის სიხშირე, სიცოცხლის
ხანგრძლივობა და განვითარების სიჩქარე მკვეთრად ეცემა, ვირუსის რეპლიკაციაც მცირდება. Watts et al.-ის ექსპერიმენტში
20°C-ზე ვირუსი კოღოს სანერწყვე ჯირკვლებამდე ვერ აღწევდა, გადაცემა კი მხოლოდ 30°C-ზე და მეტზე ხდებოდა.
\imp{ერთთვიანი lag მინიმალური ტემპერატურისთვის} სავარაუდოდ ზრდასრული კოღოს გადარჩენაზე პირდაპირ გავლენას ასახავს,
ხოლო \imp{სამთვიანი lag მაქსიმალური ტემპერატურისთვის} კოღოების პოპულაციის დინამიკაზე კუმულაციურ ეფექტს.

ნალექის როლი მოსალოდნელზე მცირე აღმოჩნდა (ერთი თვით დაყოვნებული ნალექი 54-დან მხოლოდ მე-11 ადგილზეა).
ქალაქებში ხელოვნური წყლის ჭურჭელი კოღოს გამრავლების ადგილს მთელი წლის განმავლობაში იძლევა, ნალექისგან
დამოუკიდებლად. ძლიერმა წვიმამ შეიძლება ლარვები ჩარეცხოს, გვალვის დროს კი ხალხი წყალს სახლში ინახავს,
რაც გამრავლების ახალ ადგილებს ქმნის.

\begin{etl}
ეს ორი აბზაცი გეხმარება აჩვენო, რატომ არის \textbf{მოქნილი ცვლადების არჩევა} მნიშვნელოვანი: მკვლევარმა
წინასწარ არ იცის, რომელი ცვლადი (ნალექი? Tmin? Tmax?) და რომელი lag აღმოჩნდება მნიშვნელოვანი. ამიტომ ბევრი
ვარიანტის სწრაფად მომზადება სჭირდება. შენი სისტემის „alternative data sources“ და ცვლადების
დამატების ფუნქცია ზუსტად ამას ემსახურება (ხელმძღვანელიც სწორედ ეს აღნიშნა).
\end{etl}

\textbf{სოციოეკონომიკური ფაქტორები და ურბანიზაცია} \textit{[შემოკლებით]}. GPP-ის მაღალი მნიშვნელობა
„პარადოქსულ“ კავშირს აჩვენებს: მეტი ეკონომიკური განვითარება მეტ რისკს ნიშნავს. ეს ტაილანდში დენგეს
ურბანულ ბუნებას ასახავს: ხელოვნური წყლის ჭურჭელი, მოსახლეობის მაღალი სიმჭიდროვე, ურბანული სითბოს კუნძულის ეფექტი.

\textbf{COVID-19 პანდემიის გავლენა} \textit{[შემოკლებით]}. მობილობის ცვლილებამ ადამიანისა და კოღოს კონტაქტი შეამცირა,
ხოლო ჯანდაცვაზე მიმართვის შემცირებამ შემთხვევების აღრიცხვაც შეამცირა.

\begin{citat}
\textit{„Future modeling efforts should consider incorporating real-time mobility data, healthcare utilization patterns, and other behavioral indicators to maintain predictive accuracy during unprecedented circumstances.“} (Discussion, p.~18)
\end{citat}

\begin{etl}
ეს ციტატა \imp{შენი სისტემის გაფართოებადობის (extensibility) მოტივაციაა}: ავტორები ამბობენ, რომ მომავალში
\textbf{კიდევ მეტი სხვადასხვა წყაროს} (მობილობა, ჯანდაცვა, ქცევა) გაერთიანება იქნება საჭირო. რაც მეტია წყარო,
მით მეტია ETL სამუშაო და მით უფრო სასარგებლოა ავტომატიზებული სისტემა, რომელსაც ახალი წყაროს დამატება შეუძლია.
\end{etl}

\subsection*{Limitations (შეზღუდვები): სრული თარგმანი} \pg{18}

რამდენიმე შეზღუდვა ყურადღებით უნდა იქნას გათვალისწინებული. \textbf{პირველი}, პასიურ ზედამხედველობაზე
დამოკიდებულება სავარაუდოდ დენგეს რეალური ინციდენტობის სისტემატურ შეფასებას იწვევს, განსაკუთრებით
რეგიონებში, სადაც ჯანდაცვაზე წვდომა შეზღუდულია. მხოლოდ DHF შემთხვევების გამოყენებამ შეიძლება
მძიმე ფორმებისკენ მიკერძოება გამოიწვიოს. \textbf{მეორე}, \imp{ვექტორის სიმრავლის, ადამიანების მობილობისა და
იმუნიტეტის პირდაპირი საზომების არარსებობა} გადაცემის დინამიკის გაგებას ზღუდავს. \textbf{მესამე}, დენგეს
კონტროლის ზომების ეფექტი არ იქნა გათვალისწინებული. \textbf{მეოთხე}, ბინარული კლასიფიკაცია მაღალ სიზუსტეს
შეიძლება ნაწილობრივ იმით აღწევდეს, რომ ის მუდმივად ენდემურ პროვინციებს ცნობს და არა დროით ცვლილებებს.
\textbf{მეხუთე}, ბინარული კლასიფიკაცია რისკის სურათს შეიძლება ზედმეტად ამარტივებდეს. მრავალკლასიანი მიდგომა
უფრო დეტალურ სურათს მოგვცემდა. ეს რომ შეგვემოწმებინა, ანალიზი უფრო მკაცრი, 75-ე პროცენტილის ზღვრით
(6.08 / 100\,000) გავიმეორეთ: AUC 0.79 (COVID-19-მდე) და 0.72 (პანდემიის დროს და მის შემდეგ), ტემპერატურის
დომინირება კი შენარჩუნდა. \textbf{მეექვსე}, მოდელის ხარისხის დაცემა პანდემიის დროს უჩვეულო საზოგადოებრივი
ცვლილებების მიმართ მის მოწყვლადობას აჩვენებს. \textbf{და ბოლოს}, ნარჩენების მოკლე lag-ების (1--3 თვე)
ავტოკორელაცია მნიშვნელოვანი აღმოჩნდა, რადგან ავტორებმა ავტორეგრესიული წევრები შეგნებულად არ ჩართეს.
ოპერატიული პროგნოზირებისთვის ავტორეგრესიული სტრუქტურა იქნებოდა საჭირო.

\begin{etl}
მეორე შეზღუდვა („absence of direct measures of vector abundance, human mobility, and immunity“) კიდევ ერთი
არგუმენტია: \imp{ავტორებს მეტი მონაცემი სჭირდებოდათ, ვიდრე ჰქონდათ.} შენი სისტემის წყაროების ავტომატური
მოძიება (source discovery) სწორედ ამ პრობლემას ეხება. მაგრამ ფრთხილად: ავტორები არ ამბობენ, რომ ეს
მონაცემები \textit{ვერ იპოვეს}, მხოლოდ ის, რომ \textit{არ ჰქონდათ}.
\end{etl}

%==========================================================
\section{Conclusions (დასკვნები): სრული თარგმანი} \pg{18--19}

XGBoost-ისა და SHAP-ის გამოყენებით ტაილანდში 20 წლის განმავლობაში პროვინციის დონის დენგეს რისკის
გარემო და სოციოეკონომიკური პრედიქტორები გამოვავლინეთ (ტესტის AUC = 0.80 პანდემიამდე; 0.74 პანდემიის
დროს და მის შემდეგ). ტემპერატურა დომინანტური სიგნალი აღმოჩნდა: ათი ყველაზე გავლენიანი ცვლადიდან შვიდი
ტემპერატურას ეხება, არაწრფივი ზღვრებით დაახლ. 21°C-ზე (ერთი თვით დაყოვნებული მინიმალური ტემპერატურა)
და 32°C-ზე (სამი თვით დაყოვნებული მაქსიმალური ტემპერატურა). \imp{ერთ- და სამთვიანი lag-ების სიჭარბე
მიუთითებს, რომ დენგეს რისკს ტაილანდში არა მიმდინარე პირობები, არამედ თერმული ისტორია განსაზღვრავს}.
ეს ექსპერიმენტულ ლიტერატურაში აღწერილ ვირუსის ინკუბაციის პერიოდისა და კოღოების დემოგრაფიის
ტემპერატურაზე დამოკიდებულებას შეესაბამება. პროვინციის უფრო მაღალი მთლიანი პროდუქტი გაზრდილ რისკთან იყო
დაკავშირებული, რაც დენგეს ძირითადად ურბანულ ხასიათს შეესაბამება. თუმცა, GPP მოსახლეობის სიმჭიდროვეს,
ჯანდაცვის ინფრასტრუქტურასა და აღრიცხვის მოცვასაც ასახავს და ის წმინდა ურბანიზაციის სიგნალად არ უნდა
იქნას აღქმული.

COVID-19 პერიოდმა გარემოზე დაფუძნებული პროგნოზირების მნიშვნელოვანი შეზღუდვები გამოავლინა: 2021 წელს
AUC 0.62-მდე დაეცა და მოდელი რისკს სისტემატურად აჭარბებდა. ამიტომ ვურჩევთ, რომ დიდი ეპიდემიოლოგიური
რყევების დროს მოდელის ალბათობები საზოგადოებრივი ჯანდაცვის გადაწყვეტილებებისთვის მოდელის ხელახალი
კალიბრაციის გარეშე არ იქნას გამოყენებული. უფრო ზოგადად, ჩვენი შედეგები მიუთითებს, რომ
\imp{დენგეს პროგნოზირების სისტემების შემდეგ თაობას გარემოს სიგნალების გაერთიანება დასჭირდება მობილობის,
ვექტორების პოპულაციისა და ზედამხედველობის მოცვის რეალური დროის ინდიკატორებთან}, რათა საზოგადოებრივი
ცვლილებების დროსაც საიმედო დარჩეს.

\begin{citat}
\textit{„More broadly, our findings suggest that the next generation of dengue prediction systems will need to combine environmental signals with real-time indicators of mobility, vector populations, and surveillance coverage to remain reliable through periods of societal change.“} (Conclusions, p.~19)
\end{citat}

%==========================================================
\section{შეჯამება შენთვის: სტატია vs. შენი სისტემა}

\begin{longtable}{@{}p{0.05\linewidth}p{0.42\linewidth}p{0.10\linewidth}p{0.36\linewidth}@{}}
\toprule
\textbf{\#} & \textbf{ავტორების ხელით შესრულებული ნაბიჯი} & \textbf{შენთან} & \textbf{კომენტარი}\\
\midrule
1 & დენგეს შემთხვევების მოპოვება (ეპიდ. ბიურო, პროვინცია, თვე) & $\sim$ & OpenDengue, მაგრამ მხოლოდ ქვეყნის დონეზე\\
2 & ადმინისტრაციული ერთეულების ჰარმონიზაცია (ბუენგკანი $
\rightarrow$ ნონგ ხაი) & $\times$ & ახალი მოთხოვნა ან შეზღუდვა\\
3 & კლიმატის მოპოვება (TerraClimate, 6 ცვლადი) & $\sim$ & ERA5/NASA POWER, მხოლოდ T\textsubscript{mean} და ნალექი\\
4 & გრიდის სივრცითი გასაშუალოება პროვინციაზე & $\times$ & ახლა ერთი წერტილი: ახალი მოთხოვნა\\
5 & lag 1--3 სვეტების გამოთვლა & $\sim$ & მხოლოდ კორელაცია (FR-19): Gemini-ს დავალება B\\
6 & ბიოკლიმატური ინდექსები (bio1--bio19) & $\times$ & შესაძლო Future Work\\
7 & სოციოეკონომიკური მონაცემები 3 უწყებიდან & $\sim$ & World Bank / UN WPP (FR-17), ქვეყნის დონე\\
8 & გამოტოვებული წლების წრფივი ინტერპოლაცია & ? & შეამოწმე, აკეთებს თუ არა ამას შენი სისტემა\\
9 & წლიური $
\rightarrow$ თვიური (განმეორებით) & $\checkmark$ & შენ mass-preserving მეთოდი გაქვს (FR-20); შეადარე\\
10 & ინციდენტობა 100\,000-ზე & $\checkmark$ & FR-17\\
11 & სატელიტური ინდექსები (NDVI, MNDWI) & $\times$ & სისტემის ფარგლებს გარეთაა\\
12 & ყველაფრის ერთ თვიურ ცხრილად გაერთიანება & $\checkmark$ & შენი სისტემის ბირთვი\\
13 & ნაბიჯების დოკუმენტირება (provenance) & $\checkmark$ & შენთან ავტომატურია; სტატიაში ზოგი დეტალი აკლია (მაგ. როგორ გაასაშუალოვეს)\\
\bottomrule
\end{longtable}
{\small $\checkmark$ = აქვს, $\sim$ = ნაწილობრივ, $\times$ = არ აქვს, ? = შესამოწმებელია}

\subsection*{რა გააკეთო ახლა}
\begin{enumerate}[leftmargin=1.5em]
  \item ორიგინალში წაიკითხე \textbf{Methods (გვ.~4--8)} და \textbf{Table 1 (გვ.~6--7)}. ლურჯი ციტატები ორიგინალში იპოვე და მონიშნე.
  \item გადაამოწმე, არის თუ არა \textbf{Peter Haddawy} შენი ხელმძღვანელის წერილის „Peter“.
  \item ზემოთ მოცემული ცხრილი შენი კითხვის ფურცლის (\texttt{kitxvis\_furceli.md}) მე-3 სექციაში გადაიტანე და შეავსე.
  \item Zenodo-ზე ავტორების მონაცემები ნახე (\href{https://doi.org/10.5281/zenodo.21053822}{10.5281/zenodo.21053822}): შენი სისტემის შედეგის მათთან შედარება Case Study-ს ძალიან ძლიერს გახდის.
  \item ხელმძღვანელს ცხრილის $\times$-ები აჩვენე და ჰკითხე, რომელი დაამატო სისტემაში.
\end{enumerate}

\end{document}
